\section{Limitations and Reproducibility}
\label{sec:limitations}

\paragraph{Scope.}
The completed semantic-operator experiment is a binary sentiment classifier and
the second workload is multiple-choice QA. This is insufficient to claim that
the learned policy generalizes to semantic joins, maps, extraction, free-form
generation, or multi-model pools. Each workload uses one Mini/Nano pair, and
provider updates can change outcome labels. A final VLDB submission should add
at least two operators and one open-weight answer-model pair.

\paragraph{Selection and statistical power.}
The historical 82-row Movie test was examined repeatedly while architectures,
losses, sampling mixtures, and prompt lengths were developed. It is therefore a
development test, not a clean estimator of the final chosen system. The
200-row manifest was genuinely untouched for one locked run, but became
consumed afterward. Its 16 critical cases still provide limited power, and the
40-row validation split contains only four. We disclose this chronology and do
not convert post-hoc variants into fresh claims. The next evaluation must be
pre-registered and drawn from still-unpaired reviews.

\paragraph{Acquisition assumptions.}
Disagreement-guided acquisition assumes that both candidate models can screen
the unlabeled pool and that the task supplies a fixed way to compare their
answers. It reduces human gold-annotation demand, not paired answer-model
calls. For free-form generation, semantic equivalence errors can corrupt the
agreement split; that evaluator must be fixed before selection. The policy can
also miss systematic shared failures because an agree-but-wrong pair is not
prioritized. Our benchmark experiments replay cached outputs and mask the
provided answers until selection; they estimate which human labels would have
been requested rather than measuring a new annotation campaign.

\paragraph{Objective choice.}
Inverse expected route cost removes hand-tuned outcome constants, but its
training-split estimate can shift across workloads, providers, and price
snapshots. A route-level mean also hides per-example length variation, and
inverse-cost weighting does not itself emphasize rare recoverable failures.
Our active acquisition stage must supply those critical cases. We therefore
report the full outcome matrix and will test price perturbations, alternative
cost estimators, and explicit quality constraints.

\paragraph{Cost and systems integration.}
Movie results report Mini calls, not end-to-end dollars or latency. A local
router is not free, and an answer-aware cascade always incurs a Nano call.
Before submission we will measure wall-clock routing overhead, batch
throughput, serialized prompt bytes, and realized token cost under a common
execution harness. We will also evaluate how tuple-level errors propagate to
query-level counts and ratios, which is the data-management consequence of the
operator decision.

\paragraph{Artifact plan.}
The repository already records split JSON files, locked manifests, paired
answer caches, prompt candidates, exact-request router caches, per-example
traces, and summary JSON/CSV files. The submission artifact will add a single
environment specification, immutable model/version identifiers, commands for
every table, and checks that no review ID crosses a split. It will separate
exploratory and locked outputs and include the fresh-manifest hash. Because
PVLDB Volume 20 requires an artifact link with the initial EA\&B submission,
the placeholder in the top matter must be replaced before submission.
